\documentclass[11pt]{article}

\newcommand{\cnum}{Math 245C}
\newcommand{\ced}{Spring 2026}
\newcommand{\ctitle}[4]{\title{\vspace{-0.5in}\cnum, \ced\\week #1: #2}\author{\vspace{-0.35in}\\#3, #4}}
\usepackage{enumitem}
\usepackage[usenames,dvipsnames,svgnames,table,hyperref]{xcolor}
\usepackage{amsmath, amsfonts}
\usepackage{graphicx} % Required for inserting images
\usepackage{float}
\usepackage{listings}
\usepackage{xcolor}
\usepackage{hyperref}
\usepackage{comment}

\renewcommand*{\theenumi}{\alph{enumi}}
\renewcommand*\labelenumi{(\theenumi)}
\renewcommand*{\theenumii}{\roman{enumii}}
\renewcommand*\labelenumii{\theenumii.}
\newcommand{\notiff}{%
  \mathrel{{\ooalign{\hidewidth$\not\phantom{"}$\hidewidth\cr$\iff$}}}}

\definecolor{codegreen}{rgb}{0,0.6,0}
\definecolor{codegray}{rgb}{0.5,0.5,0.5}
\definecolor{codepurple}{rgb}{0.58,0,0.82}
\definecolor{backcolour}{rgb}{0.95,0.95,0.92}
\definecolor{header}{HTML}{205459}
\definecolor{solution}{HTML}{204d52}

\newcommand{\solution}[1]{{{\textcolor{header}{Solution:} \textcolor{solution}{#1}}}}
\setlist[enumerate]{label=(\alph*)}

\lstdefinestyle{mystyle}{
    backgroundcolor=\color{backcolour},
    commentstyle=\color{codegreen},
    keywordstyle=\color{magenta},
    numberstyle=\tiny\color{codegray},
    stringstyle=\color{codepurple},
    basicstyle=\ttfamily\footnotesize,
    breakatwhitespace=false,
    breaklines=true,
    captionpos=b,
    keepspaces=true,
    numbers=left,
    numbersep=5pt,
    showspaces=false,
    showstringspaces=false,
    showtabs=false,
    tabsize=2
}


\begin{document}
\ctitle{\#1}{}{Asher Christian}{}
\date{}
\maketitle

\section{Monday}
Theorem: Let $f : \mathbb{R}^{d} \to \mathbb{C}$ be a  locally  bounded Borel function such that
$|f| \equiv 1$ and $f(x+y) = f(x)f(y)$ for all $x,y \in \mathbb{R}^{d}$. Then
\begin{enumerate}
    \item there exists $\xi \in \mathbb{R}^{d}$ such that
        \[
        f(x) = e^{2\pi i \langle x, \xi \rangle}
        \] 
        for $x \in \mathbb{R}^{d}$ 
    \item If we further assume that $f : \Pi^{d} = \mathbb{R}^{d} / \mathbb{Z}^{d} \to \mathbb{C}$ 
        then $\xi \in \mathbb{Z}^{d}$ $f(x+ k) = f(x) $ for all  $x \in \mathbb{R}^{d}$ and $k \in \mathbb{Z}^{d}$
\end{enumerate}
Proof:
Let $\{e_1,\dots,e_d\}$  be the standard orthonormal
basis of $ \mathbb{R}^{d}$. Set
\[
f_j(t) = f(te_j) \quad \forall t \in \mathbb{R}
\] 
by Friday exercise, there exists $\xi_j \in \mathbb{R}$ such that
\[
f_j(t) = e^{2\pi i t \xi_j} \quad \forall t \in \mathbb{R}
\] 
Note that
\[
    f(x) = f(\sum_{j=1}^{n}x_j e_j) = \Pi_{j=1}^{n} f(x_je_j) = \Pi_{j=1}^{n}f_j(x_j)
\] 
\[
 = e^{2\pi i \langle x, \xi \rangle}
\] 
where $\xi = (\xi_1,\dots,\xi_d)$
Goal: Show that $(f_{\xi})_{\xi \in \mathbb{Z}^{d}}$ is an orthonormal basis of $ L^{2}( \Pi^{d},  \mathbb{C})$
\[
    L^{2}( \Pi^{d}, \mathbb{C}) = \bigcup_{k=1}^{\infty} span\{f_{\xi_1},\dots,f_{\xi_k}\} \quad \{f_{\xi}\}_{\xi \in \mathbb{Z}^{d}} = \{f_{\xi_j}\}_{j \in \mathbb{N}}
\] 
The weierstrass approximation theorem tells us that $\overline{p[0,1]^{d}}^{C([0,1]^{d})} = C([0,1]^{d})$ 
Stone-weierstrass theorem

$S, +, \cdot$ is an algebra if it is a vector space with respect to addition and a ring with respect to multiplication\\
Let $ \mathcal{X}$ be locally compact hausdorff space
Step 1: $ \mathcal{X}$ is a compact Hausdorff space
Step 2: $ \mathcal{X}$ is a locally compact hausdorff space

Let $S \subset C_0( \mathcal{X}, \mathbb{R}) = C_0( \mathcal{X})$ we call $ S$ a sub-algebra of  $ C_{0} ( \mathcal{X})$ if
$S$ is a real vector space such that
 \[
fg \in S \impliedby f,g \in S
\] 
we call $S$ a lattice if
 \[
     \min\{f,g\}, \max\{f,g\} \in S \impliedby f,g \in S
\] 
We say that $S$ separates points if
whenever  $x,y \in X$ and  $x \ne y$ there exists  $f \in S$ such that
 \[
f(x) \ne f(y)
\] 
we say that $S$ vanishes at  $x_0 \in X$ if
\[
\forall f \in S, \quad f(x_0) = 0
\] 
Remark: If we replace $ C_0( \mathcal{X})$ by  $ C_0( \mathcal{X}, \mathbb{C})$
and the above definition (except for lattice) can be extended
\section{Theorem: Stone-Weierstrass, Real Case}
Here and in next theorem we further assume that
$ \mathcal{X}$ is compact. $C_0( \mathcal{X}) = C( \mathcal{X}) = C_c( \mathcal{X})$.
Let $ A$ be  a sub-algebra  of $ C_0( \mathcal{X})$ then the following are equivalent
\begin{enumerate}
    \item 
        \[
        \overline{ A}^{C_0( \mathcal{X})} = C_0( \mathcal{X})
        \] 
    \item $A$ separates points and vanishes nowhere
\end{enumerate}
Remark [compact]: Assume  $A$ is a closed sub-algebra of   $C_0 ( \mathcal{X})$ then either 
$A = C_0( \mathcal{X})$ or there exists $x_0 \in \mathcal{X}$ such that
\[
    A = \overline{A}^{C_0( \mathcal{X})} = \{f \in C_0( \mathcal{X}) : f(x_{0}) = 0\}
\] 
Corollary: If $  \mathcal{X} \subset \mathbb{R}^{d}$ is compact, and $ \mathbb{R}[x]$ is the set of polynomials
on $ \mathcal{X}$ then $ C( \mathcal{X}) = \overline{ \mathbb{R}[ \mathcal{X}]}^{C( \mathcal{X})}$\\
Proof:
\[
    \mathbb{R}[ \mathcal{X}]
\] 
separates points in $ \mathcal{X}$ and is nowhere vanishing
%But polynomials are not dense in $C_0( \mathbb{R})$ right?

\section{Theorem: Stone-Weierstrass, The complex case}
Let $A$ be a closed complex sub-algebra of  $C_{0}( \mathcal{X}, \mathbb{C})$ which is closed under conjugation
and separates points, then one of the following holds.
\begin{enumerate}
    \item $A = C_0( \mathcal{X}, \mathbb{C})$
    \item $\exists x_0 \in \mathcal{X}$ such that $A = \{f \in C_0( \mathcal{X}, \mathbb{C}) : f(x_0) = 0\}$
\end{enumerate}
Proof:
\[
f = f_1 + if_2 \in A \implies \overline{f} = f_1 - if_2 \in A \implies f_1 \in A
\] 
implies that
\[
if = if_1 - f_2 \in A \implies f_2 \in A
\] 
Let 
\[
    A_{ \mathbb{R}} = \{R(f) : f \in A\}
\] 
we have that
\[
    A = A_{ \mathbb{R}} + i A_{ \mathbb{R}}
\] 
We have thatmk $ A_{ \mathbb{R}}$ is a closed sub-algebra of $ C_0 ( \mathcal{X})$ By the previous remark,
either $ A_{ \mathbb{R}} = C_0( \mathcal{X})$  or $A_ \mathbb{R}$ vanishes at some point
If $A_{ \mathbb{R}} = C_0( \mathcal{X})$ then
\[
A = C_0( \mathbb{R}) + i  C_0( \mathbb{R}) \implies A = C_0( \mathcal{X}, \mathbb{C})
\] 
this concludes the proof.

\section{Wednesday}
Fourier Series and Transform:\\
Reminder: If $X$ is a compact Hausdorff space and  $A \subset C( \mathcal{X}, \mathbb{C})$ is a complex
closed sub-algebra closed under conjugation and separates points.
Then either $A = C( \mathcal{X}, \mathbb{C})$ or else there exists $x_0 \in \mathcal{X}$
such that
\[
    A = \{f \in C( \mathcal{X}, \mathbb{C}) : f(x_0) = 0\}
\]  %isnt i just a subset?
Warning:
\[
    p([0,1]^{d}, \mathbb{C}) :=    C([0,1]^{d}, \mathbb{C}) \cap \mathbb{C}[x_1,\dots,x_d]
\] 
%really? 
is not closed under conjugation. In fact $p([0,1]^{d}, \mathbb{C})$ is contained in the set of holomorphic
function on $(0,1)^{d}$ \\
\[
    \mathbb{R}[x_1,\dots,x_d]
\] 
polynomail of real coefficients
\[
 \mathbb{R} \langle x_1,\dots,x_d \rangle
\]
for $x_1,\dots,x_d \in A$ an algebra for example $\Pi_n ( \mathbb{R})$ the set of $n \times n$ matrices.
\\
Notation
 \[
     \Pi^{d} = \mathbb{R}^{d} / \mathbb{Z}^{d} \ni [x] = \{x + k : k \in \mathbb{Z}^{d}\}
\] 
\begin{enumerate}
    \item For $k_1,\dots,k_d \in \mathbb{Z}$ we define
        \[
        E_k(x) = e^{2\pi  i \langle x, k \rangle} \quad x \in \mathbb{R}^{d}
        \] 
        where $\langle x , k \rangle = \sum_{j=1}^{d}x_jk_j$
    \item If $p,q \in [1,+\infty]$ and $\frac{1}{p} + \frac{1}{q} = 1$ for $\Omega \subset \mathbb{R}^{d}$ 
        Borel and $f \in L^{p}(\Omega, \mathbb{C})$ $g \in L^{q}(\Omega, \mathbb{C})$
        we set
        \[
        \overline{g} = \Re(g) + i \Im(g)
        \] 
        \[
            \langle f, g \rangle_{\Omega} := \int_{ \Omega}^{}f(x) \overline{g}(x) dx
        \] 
\end{enumerate}
Exercise: Let $f \in L^{1}_{loc}( \mathbb{R}^{d}, \mathbb{C})$ be such that
 \[
f : \Pi^{d} \to \mathbb{C}
\] 
(is periodic). Show that for any  $\xi \in \mathbb{R}^{d}$ 
\[
    \int_{[0,1)^{d}}^{}f(x)dx = \int_{[0,1)^{d} + \xi}^{}f(x)dx
\] 
Assume without loss of generality that $d = 1$ and  $\xi \in [0,1)$
we are to show that
 \[
     \int_{\xi}^{xi + 1}f(t)dt = \int_{0}^{1}f(t)dt
\] 
If we further assume that $f$ is continuous then 
\[
\xi \to \int_{\xi}^{\xi + 1}f(t)dt = a(\xi)
\] 
is differentiable and
\[
a'(\xi) = f(\xi+1) - f(\xi) = 0
\] 
and so $a(\xi) = a(0)$
Let  $\rho \in C_c^{\infty}([-1,1])$ be a probability density 
and set $\rho_\epsilon(t) = \frac{1}{\epsilon}\rho(\frac{t}{\epsilon})$ then
\[
    \rho_{\epsilon} * f \in C^{1}(\Pi)
\] 
and so
\[
    \int_{\xi}^{\xi+1}f(t)dt = \lim_{\epsilon\to 0} \int_{\xi}^{\xi+1} \rho_{\epsilon} * f(t) = \lim_{\epsilon\to 0 } \int_{0}^{1}\rho_{\epsilon} * f(t)dt = \int_{0}^{1}f(t)dt
\] 
Exercise:
Show that the set
\[
\{E_k : k \in \mathbb{Z}\}
\] 
is an orthonormal basis for $L^{2}(\Pi^{d}, \mathbb{C})$ \\
Solution:
\[
    \overline{C( \Pi^{d}. \mathbb{C})}^{L^2(\Pi^{d}, \mathbb{C})} = L^2(\Pi^{d}, \mathbb{C})
\] 
If we can show
\[
\overline{S}^{C(\Pi^{d}, \mathbb{C})} = C(\Pi^{d}, \mathbb{C})
\] 
Note that if $k, l \in \mathbb{Z}^{d} $ then
\[
    \langle E_k, E_l \rangle_{L^2(\Pi^{d}, \mathbb{C})} = \int_{\Pi^{d}}^{}e^{2\pi i \langle k-l,x \rangle}dx
\] 
\[
 = \Pi_{j=1}^{d} \int_{0}^{1} e^{2\pi i (k_j - l_j)x_j}dx_j = 
\begin{cases}
    1 & k = l\\
    0 & k \ne l
\end{cases}
\] 
Let $S = \overline{span\{ E_k : k \in \mathbb{Z}^{d}\}}^{C(\Pi^{d}, \mathbb{C})}$ 
The set $S$ is a closed sub-algebra of  $C(\Pi^{d}, \mathbb{C})$ which separates points and is closed
under conjugation and since $E_{0,\dots,0} \in S$ we conclude that
$S = C( \Pi^{d}, \mathbb{C})$\\
Definition:
If $f \in L^2( \Pi^{d})$ and $g \in L^{1}( \mathbb{R}^{d})$ for
$k \in \mathbb{Z}^{d}$ and $\xi \in \mathbb{R}^{d}$ we define
\[
    \hat f(k) = \langle f, E_k \rangle_{L^2(\Pi^{d}, \mathbb{C})} 
\] 
\[
    F(g)(\xi) = \int_{ \mathbb{R}^{d}}^{}g(x)e^{-2\pi i \langle \xi ,x\rangle}dx= \langle g, \xi \rangle_{ \mathbb{R}^{d}}
\] 
we call
\[
    \hat f
\] 
the fourier transform of $f$ and
 \[
F(g) 
\] 
the fourier transform of $g$.\\
Fact: 
\[
    f =^{||\cdot||_{L^2}} \sum_{k \in \mathbb{Z}^{d}}^{} \hat f(k) E_k
\] 
We call the sum
\[
\sum_{k \in \mathbb{Z}^{d}}^{} \hat f(k) E_k
\] 
the fourier series of $f$
Proof: Define  $f_N = \sum_{|k| \le N}^{} \hat f(k) E_k$  this is the projection
\[
    proj_{S_N} f
\] 
where
\[
    S_n = span\{E_k : |k| \le N\}
\] 

\section{Friday}
If $k \in \mathbb{Z}^{d}$ $E_k(x) = e^{2\pi i \langle k , x \rangle}$ for all $x \in \Pi^{d}$
\begin{enumerate}
    \item $(E_k)_{k \in \mathbb{Z}^{d}}$ is an orthonormal basis 
        for $L^2(\Pi^{d}, \mathbb{C})$
\end{enumerate}
Remark: Let $f \in L^2(\Pi^{d}, \mathbb{C})$. Then
\[
f = \sum_{k \in \mathbb{Z}^{d}}^{}\hat f(k) E_k \quad \hat f(k) = \langle f, E_k \rangle
\] 
Note that $\hat f : \mathbb{Z}^{d} \to \mathbb{C}$ and
\[
    (*) \quad   ||f||_{L^2(\Pi^{d}, \mathbb{C})}^{2} = \sum_{k \in \mathbb{Z}^{d}}^{}|\hat f (k)|^2
\] 
\[
    = ||\hat f||_{L^2( \mathbb{Z}^{d}, \mathcal{H}^{0})}^{2}
\] 
where, here
\[
    \mathcal{H}^{0}
\] 
is the counting measure on $ \mathbb{Z}^{d}$ 
\[
    L^2( \mathbb{Z}^{d}, \mathcal{H}^{0}) = l_2( \mathbb{C})
\] 
If $f \in L^{1}(\Pi^{d}, \mathbb{C})$ then $|\hat f(k)| \le ||f||_{L^{1}(\Pi^{d}, \mathbb{C})}$ since $||E_k||_{L^{\infty}(\Pi^{d}, \mathbb{C})} = 1$
Hence: 
\[
    (**) \quad    ||\hat f||_{l_{\infty}( \mathbb{C})} = ||\hat f||_{L^{\infty}(Z^{d}, \mathcal{H}^{0})} \le ||f||_{L^{1}( \Pi^{d}, \mathbb{C})}
\] 
then $(*)$ is Parseval's identity
then
\[
\begin{cases}
    L^{2}(\Pi^{d}, \mathbb{C}, \mathcal{L}^2) & L^2( \mathbb{Z}^{d}, \mathbb{C}, \mathcal{H}^{0})\\
    L^{1}(\Pi^{d}, \mathbb{C}, \mathcal{L}^2) & L^{\infty}( \mathbb{Z}^{d}, \mathbb{C}, \mathcal{H}^{0})
\end{cases}
\] 
Can we conclude that for $1 < p < 2$
\[
    L^{p}(\Pi^{d}, \mathbb{C}, \mathcal{L}^{2}) \to L^{q}( \mathbb{Z}^{d}, \mathbb{C}, \mathcal{H}^{0})
\]?\\
Goal: let $(\Omega_0, \mu_0)$ be a $\sigma$-finite measure space
and let $(\Omega_1, \mu_1)$ be another $\sigma$-finite measure space. Assume that
\[
    1 \le p_0 < p_1 \le +\infty
\] 
\[
1 \le q_0, q_1 < +\infty
\] 
Define the interpolation $p_t$ and  $q_t$ by
 \[
\frac{1}{p_t} = \frac{1-t}{p_0} + \frac{t}{p_1} \qquad \frac{1}{q_t} = \frac{1-t}{q_0} + \frac{t}{q_1}
\] 
Assume that
\[
T_0 : L^{p_0}(\Omega_0, \mathbb{C}, \mu_0) \to L^{q_0}(\Omega_1, \mathbb{C}, \mu_1)
\] 
linear, bounded, and 
\[
T_1: L^{p_1}( \Omega_0, \mathbb{C}, \mu) \to L^{q_1}(\Omega_1, \mathbb{C}, \mu_1)
\] 
is linear and bounded. We want to show that $T$ yields
\[
T_t : L^{q_t}(\Omega_0, \mathbb{C}, \mu_0) \to L^{q_t}(\Omega_1, \mathbb{C}, \mu_1)
\] 
is linear bounded
\[
    (***) \quad    ||T_t||_{L^{p_t} \to L^{q_t}} \le ||T_0||_{L^{p_0}\to L^{q_0}}^{1-t}||T_1||_{L^{p_1} \to L^{q_1}}^{t}
\] 
\[
\log(\dots) \le (1-t)\log + t\log
\] 
We show that $(***)$ is a consequence of the following
we draw the complex plan  $ \mathbb{C} = \mathbb{R} + i \mathbb{R} \sim \mathbb{R}^2$.
Facts from complex analysis:
\[
    f: \mathbb{C} \to \mathbb{C} \quad \text{holomorphic}
\] 
\[
    M_0 = \sup\{|f(x)| : \Re(z) = 1\} \qquad M_t = \sup\{|f(z)| : \Re(z) = (1-t)1 + t2\}
\] 
Then
\[
M_t \le M_0^{1-t}M_1^{t}
\] 
Elements of complex analysis.
We identify $ \mathbb{C}$ with $ \mathbb{R}^2$ where $x +iy \to(x,y)$
and similarly we identify  $D \subset \mathbb{C}$ by $\tilde D \subset \mathbb{R}^2$
and $f : D \to \mathbb{C}$ with $\tilde f : \tilde  D \to \mathbb{C}$
Let
\[
D \subset \mathbb{C}
\] 
be an open set.
We say that $f : D \to \mathbb{C}$ is holomorphic if
\[
f'(z) 
\] 
exists for every $z \in D$.
We say that  $u : D \to \mathbb{R}$ is harmonic if
\[
u \in C^2( D)
\] 
and
\[
\frac{\partial^2 u}{ \partial x^2} + \frac{\partial^2 u}{ \partial y^2} = \triangle u \equiv 0
\] 
holomorphic functions have real and imaginary part harmonic.
Fact: Assume that $u \in C^2(D)$ is harmonic
then if $B_r(x_0) \subset D$ thena
\[
    (1) \quad   f(x_0) = \frac{1}{ \mathcal{L}^{s}(B_r(x_0))}\int_{\partial B_r(x_0)}^{} u d \mathcal{H}^{1} =  \frac{1}{ \mathcal{L}^{2}(B_r(x_0))} \int_{B_r(x_0)}^{}u d \mathcal{L}^2
\] 
Let
\[
    / \int_{A}^{}d\mu = \frac{1}{\mu[A]} \int_{A}^{}d \mu
\] 
Proof of (1) fix  $R > 0$ such that  $B_R(x_0) \subset D$
define
\[
    \phi(r) = / \int_{\partial B_r(x_0)}^{} u d \mathcal{H}^{1} = / \int_{ \partial B_1(0)}^{}u(x_0 + rx) \mathcal{H}^{1}(dz)
\] 
we have
\[
    \phi'(r) = / \int_{\partial B_1(0)}^{} \langle \nabla u(x_0 + rx), \frac{x_0 + rx - x_0}{r} \rangle \mathcal{H}^{1}(dx)
\] 
\[
    = / \int_{B_r(x_0)}^{} \langle \nabla u(x), \nu \rangle \mathcal{H}^{1}(dx)
\] 
using Gauss-Green we conclude this is equal to 
\[
    \frac{1}{\mathcal{H}^{1}(\partial B_r(x_0)} \int_{B_r(x_0)}^{} div(\nabla u) dx = 0
\] 
so $\phi : (0, R) \to \mathbb{R}$ is constant and so 
\[
    (1) \quad \phi(r) = \lim_{t\to 0}\phi(t) = u(x_0)
\] 
Note that
\[
    \int_{B_r(x_0)}^{}u(x)dx = \int_{0}^{r}ds \int_{\partial B_s(x_0)}^{}u d \mathcal{H}^{1} = \int_{0}^{r}ds \mathcal{H}^{1}(\partial B_s(x_0)) u(x_0) d \mathcal{H}^{1}
\] 
\[
    = u(x_0) \int_{ 0}^{r}\int_{\partial B_s(x_0)}^{}d\mathcal{H}^{1}ds =  u(x_0) \mathcal{L}^{d}(B_r(x_{0}))
\] 
Corollary:
If $u$ achieves its maximum on  $D$ then it achieves its maximum on  $\partial D$.
Provided that  $D$ is connected
Proof:
Let 
\[
    E = \{x \in D: u(x) = \max\}
\] 
then $E$ is closed and by 1  $E$ is open

\end{document}
